% ============================================================
%  Rapport de Stage de Fin d'Etudes
%  Abdelhafid Ibn Mhand -- DBT SDSI -- 2025/2026
%  Plateforme de Data Intelligence pour l'Analyse Statistique Territoriale
%  Projet parallele : Simpute (PyPI)
% ============================================================
\documentclass[12pt,a4paper]{report}

% ── Encoding & Language ─────────────────────────────────────
\usepackage[utf8]{inputenc}
\usepackage[T1]{fontenc}
\usepackage[provide=*,french]{babel}

% ── Arial Font via Helvetica clone ──────────────────────────
\usepackage[scaled]{helvet}
\renewcommand{\familydefault}{\sfdefault}

% ── Typography ──────────────────────────────────────────────
\usepackage[expansion=false,protrusion=true]{microtype}
\usepackage{setspace}
\onehalfspacing
\emergencystretch=3em

% ── Page Layout ─────────────────────────────────────────────
\usepackage[
  top=2.5cm, bottom=2.8cm,
  left=3cm,  right=2.5cm,
  headheight=16pt
]{geometry}

% ── Colors ──────────────────────────────────────────────────
\usepackage[dvipsnames,svgnames,x11names]{xcolor}
\definecolor{hcpblue}{RGB}{0,84,166}
\definecolor{hcpdark}{RGB}{15,20,40}
\definecolor{hcpgray}{RGB}{95,100,115}
\definecolor{hcplight}{RGB}{238,244,252}
\definecolor{hcpaccent}{RGB}{0,130,115}
\definecolor{codered}{RGB}{175,30,30}
\definecolor{covergray}{RGB}{245,247,250}
\definecolor{rulecol}{RGB}{0,84,166}

% ── Graphics & Figures ──────────────────────────────────────
\usepackage{graphicx}
\usepackage{float}
\usepackage{caption}
\usepackage{subcaption}
% Fixed figure width for consistency regardless of image resolution
\newlength{\figwidth}
\setlength{\figwidth}{0.88\textwidth}
\captionsetup{
  font={small,sf},
  labelfont={bf,color=hcpblue},
  format=hang,
  justification=raggedright,
  singlelinecheck=false,
  skip=6pt
}
\captionsetup[list=figure]{skip=4pt}

% ── Tables ──────────────────────────────────────────────────
\usepackage{booktabs}
\usepackage{tabularx}
\usepackage{multirow}
\usepackage{array}
\usepackage{longtable}
\newcolumntype{L}[1]{>{\raggedright\arraybackslash}p{#1}}
\newcolumntype{C}[1]{>{\centering\arraybackslash}p{#1}}
\newcolumntype{R}[1]{>{\raggedleft\arraybackslash}p{#1}}

% ── Code Listings (ASCII-safe) ──────────────────────────────
% Key fix: extendedchars=false prevents UTF-8 crashes inside lstlisting
\usepackage{listings}
\lstdefinestyle{codestyle}{
  basicstyle=\ttfamily\small,
  keywordstyle=\color{hcpblue}\bfseries,
  commentstyle=\color{hcpgray}\itshape,
  stringstyle=\color{hcpaccent},
  numberstyle=\tiny\color{hcpgray},
  numbers=left, numbersep=8pt,
  breaklines=true, breakatwhitespace=true,
  frame=single, framerule=0.4pt,
  rulecolor=\color{hcpblue!40},
  backgroundcolor=\color{hcplight},
  showstringspaces=false,
  tabsize=4,
  xleftmargin=14pt,
  extendedchars=false,
  columns=fullflexible,
  keepspaces=true,
}
\lstdefinestyle{treenostyle}{
  style=codestyle,
  basicstyle=\ttfamily\footnotesize,
}
\lstset{style=codestyle}

% ── Links & PDF Metadata ────────────────────────────────────
\usepackage[
  colorlinks=true,
  linkcolor=hcpblue,
  urlcolor=hcpaccent,
  citecolor=hcpblue,
  pdftitle={Plateforme de Data Intelligence pour l'Analyse Statistique Territoriale},
  pdfauthor={Abdelhafid Ibn Mhand},
  pdfsubject={Rapport de Stage -- DBT SDSI 2025/2026}
]{hyperref}

% ── Headers & Footers ───────────────────────────────────────
\usepackage{fancyhdr}
\pagestyle{fancy}
\fancyhf{}
\fancyhead[L]{\small\color{hcpgray}\textsf{Abdelhafid Ibn Mhand}}
\fancyhead[R]{\small\color{hcpgray}\textsf{Plateforme de Data Intelligence -- HCP}}
\fancyfoot[C]{\small\color{hcpgray}\thepage}
\renewcommand{\headrulewidth}{0.4pt}
\renewcommand{\footrulewidth}{0pt}
\renewcommand{\headrule}{{\color{hcpblue!50}\hrule width\headwidth height\headrulewidth}}

% ── Chapter / Section Styling ────────────────────────────────
\usepackage{titlesec}
\titleformat{\chapter}[display]
  {\sffamily\huge\bfseries\color{hcpblue}}
  {\sffamily\large\color{hcpgray}\MakeUppercase{\chaptertitlename}\ \thechapter}
  {12pt}{\Huge\sffamily\bfseries\color{hcpblue}}
\titleformat{\section}
  {\sffamily\Large\bfseries\color{hcpdark}}{\thesection}{1em}{}
\titleformat{\subsection}
  {\sffamily\large\bfseries\color{hcpblue!85!black}}{\thesubsection}{1em}{}
\titleformat{\subsubsection}
  {\sffamily\normalsize\bfseries\color{hcpgray}}{\thesubsubsection}{1em}{}
\titlespacing*{\chapter}{0pt}{8pt}{28pt}
\titlespacing*{\section}{0pt}{18pt}{8pt}
\titlespacing*{\subsection}{0pt}{14pt}{6pt}

% ── Boxes / Callouts ────────────────────────────────────────
\usepackage{tcolorbox}
\tcbuselibrary{skins,breakable}
\newtcolorbox{infobox}[1][]{
  colback=hcplight, colframe=hcpblue,
  fonttitle=\sffamily\bfseries\color{white},
  coltitle=white, colbacktitle=hcpblue,
  title={#1}, arc=4pt,
  left=8pt, right=8pt, top=5pt, bottom=5pt,
  breakable
}
\newtcolorbox{resultbox}{
  colback=hcpaccent!7, colframe=hcpaccent,
  arc=4pt, left=8pt, right=8pt, top=5pt, bottom=5pt,
  fontupper=\small\sffamily, breakable
}
\newtcolorbox{warningbox}{
  colback=codered!5, colframe=codered!65,
  arc=4pt, left=8pt, right=8pt, top=5pt, bottom=5pt,
  fontupper=\small\sffamily, breakable
}
\newtcolorbox{toolbox}[1][]{
  colback=covergray, colframe=hcpgray!60,
  fonttitle=\sffamily\bfseries\small,
  coltitle=hcpdark, colbacktitle=covergray,
  title={#1}, arc=3pt,
  left=6pt, right=6pt, top=4pt, bottom=4pt,
  breakable
}

% ── Figure placeholder (fixed size) ─────────────────────────
% Usage: \figplaceholder[short LOF caption]{filename}{caption text}{label}
\newcommand{\figplaceholder}[4][]{%
  \begin{figure}[H]
    \centering
    \fbox{%
      \begin{minipage}[c][5.5cm][c]{\figwidth}
        \centering
        {\color{hcpblue!60}\rule{\figwidth}{0.4pt}}\\[6pt]
        {\color{hcpblue}\large\bfseries\sffamily [FIGURE]}\\[4pt]
        {\color{hcpdark}\small\ttfamily #2}\\[4pt]
        {\color{hcpgray}\footnotesize\sffamily #3}\\[6pt]
        {\color{hcpblue!60}\rule{\figwidth}{0.4pt}}
      \end{minipage}%
    }
    \ifx\relax#1\relax
      \caption{#3}%
    \else
      \caption[#1]{#3}%
    \fi
    \label{#4}
  \end{figure}%
}

% When you have the real image, replace \figplaceholder[short LOF cap]{file}{cap}{lbl}
% with the standard block below (same fixed width \figwidth):
%   \begin{figure}[H]
%     \centering
%     \includegraphics[width=\figwidth]{file}
%     \caption{cap}
%     \label{lbl}
%   \end{figure}

% ── Misc ────────────────────────────────────────────────────
\usepackage{enumitem}
\usepackage{amsmath,amssymb}
\usepackage{pifont}
\usepackage{mdframed}
\usepackage{epigraph}
\setlength{\epigraphwidth}{0.65\textwidth}
\setlength{\epigraphrule}{0pt}

% Tighter list spacing
\setlist{itemsep=3pt, topsep=5pt, parsep=0pt}

% ─────────────────────────────────────────────────────────────
\hypersetup{pageanchor=false}
\begin{document}

% ============================================================
%  PAGE DE GARDE  (single page, balanced density)
% ============================================================
\begin{titlepage}
\thispagestyle{empty}

\vspace*{0.25cm}

% Logos and institution (three columns) - Real images, bigger sides, smaller text
\noindent
\begin{minipage}[c]{0.32\textwidth}
  \centering
  % Render the EST logo, maximizing the available space while preserving aspect ratio
  \includegraphics[width=\textwidth, height=2.6cm, keepaspectratio]{est_logo.png}
\end{minipage}%
\begin{minipage}[c]{0.36\textwidth}
  \centering
  \setstretch{0.85} % Tighten line spacing for smaller text layout
  {\sffamily\scriptsize\bfseries\color{hcpdark} Royaume du Maroc}\\[2pt]
  {\sffamily\tiny\color{hcpgray}
    Minist\`ere de l'Enseignement Sup\'erieur,\\
    de la Recherche Scientifique\\ et de l'Innovation}\\[3pt]
  {\sffamily\scriptsize\bfseries\color{hcpblue} Universit\'e Ibn Zohr}\\[1pt]
  {\sffamily\tiny\color{hcpgray}
    \'Ecole Sup\'erieure de Technologie\\ de Guelmim}
\end{minipage}%
\begin{minipage}[c]{0.32\textwidth}
  \centering
  % Render the HCP logo, maximizing the available space while preserving aspect ratio
  \includegraphics[width=\textwidth, height=2.6cm, keepaspectratio]{logo_hcp.png}
\end{minipage}

\vspace{0.55cm}
{\color{hcpblue}\noindent\rule{\textwidth}{0.7pt}}

\begin{center}
  {\sffamily\Large\bfseries\color{hcpgray} Rapport de Stage de Fin d'\'Etudes}\\[5pt]
  {\sffamily\normalsize\color{hcpgray}
    Organisme d'accueil~: \textbf{Haut-Commissariat au Plan,\\
    Direction R\'egionale de Guelmim}}
\end{center}

\begin{center}
  \begin{tcolorbox}[
    enhanced,
    colback=hcplight,
    colframe=hcpblue,
    arc=5pt,
    boxrule=1pt,
    left=14pt, right=14pt, top=12pt, bottom=12pt,
    width=0.92\textwidth,
    halign=center,
    drop fuzzy shadow
  ]
    {\sffamily\LARGE\bfseries\color{hcpblue}
      Plateforme de Data Intelligence}\\[6pt]
    {\sffamily\Large\bfseries\color{hcpdark}
      Analyse et Valorisation\\
      des Donn\'ees Statistiques Territoriales}\\[8pt]
    {\color{hcpblue!40}\rule{0.62\linewidth}{0.45pt}}\\[6pt]
    {\sffamily\normalsize\color{hcpdark}
      De l'ingestion des donn\'ees brutes \`a la mod\'elisation pr\'edictive\\
      et au d\'eploiement d'un tableau de bord interactif}\\[6pt]
    {\sffamily\footnotesize\color{hcpaccent}
      Projet parall\`ele~: \textbf{Simpute}, biblioth\`eque Python open source\\
      d'imputation adaptative des donn\'ees tabulaires, publi\'ee sur PyPI}
  \end{tcolorbox}
\end{center}

\begin{center}
  {\sffamily\small\color{hcpgray} Pr\'esent\'e et soutenu par}\\[5pt]
  {\sffamily\LARGE\bfseries\color{hcpdark} ABDELHAFID IBN MHAND}\\[8pt]
  {\sffamily\small\color{hcpgray} Pour l'obtention du dipl\^ome}\\[4pt]
  {\sffamily\large\bfseries\color{hcpblue}
    Dipl\^ome Bachelor Universitaire de Technologie}\\[3pt]
  {\sffamily\normalsize\color{hcpgray}
    Sp\'ecialit\'e~: Science de Donn\'ees et Syst\`emes Intelligents}
\end{center}

{\color{hcpblue!45}\noindent\rule{\textwidth}{0.4pt}}
\begin{center}
  \begin{tabular}{@{}c@{}}
    % Stacked supervisors using standard LaTeX line breaks
    {\sffamily\small\color{hcpgray} Encadrant acad\'emique~: \textbf{\color{hcpdark}Prof.\ Adnane Elmrabty}}\\[1pt]
    {\sffamily\small\color{hcpgray} Encadrant professionnel~: \textbf{\color{hcpdark}M.\ Yassine Houssa}}\\[10pt]
    
    {\sffamily\small\bfseries\color{hcpgray}
      Soutenu en Juillet 2026 devant le jury}\\[6pt]
    {\sffamily\footnotesize\color{hcpgray}
      D\'ep\^ots GitHub~:
      \href{https://github.com/Hvllvix/HCP-Internship}{\texttt{HCP-Internship}}
      \enspace$\cdot$\enspace
      \href{https://github.com/Hvllvix/Simpute}{\texttt{Simpute}}
      \enspace$\cdot$\enspace
      \href{https://github.com/Hvllvix}{\texttt{github.com/Hvllvix}}}\\[15pt]
    {\sffamily\LARGE\bfseries\color{hcpblue} Ann\'ee Universitaire 2025/2026}
  \end{tabular}
\end{center}

\end{titlepage}

% ============================================================
%  PAGES LIMINAIRES
% ============================================================
\pagenumbering{roman}
\setcounter{page}{1}

% ── Remerciements ───────────────────────────────────────────
\chapter*{Remerciements}
\addcontentsline{toc}{chapter}{Remerciements}

Avant de présenter ce travail, je tiens à exprimer ma sincère gratitude envers toutes les personnes qui ont contribué, de près ou de loin, à la réalisation de ce projet de fin d'études.

Je remercie en premier lieu la \textbf{Direction Régionale du Haut-Commissariat au Plan de Guelmim} pour m'avoir accueilli dans le cadre de ce stage. Cette expérience m'a permis de découvrir le fonctionnement d'une institution statistique publique de premier plan, de travailler directement sur des données officielles et de mesurer concrètement les enjeux de la valorisation de l'information territoriale.

J'adresse mes sincères remerciements à mon \textbf{encadrant professionnel, M.\ Yassine Houssa}, pour son suivi rigoureux, ses orientations techniques précises et sa disponibilité constante tout au long de la période de stage. Sa connaissance approfondie des sources de données du HCP et sa vision du terrain ont été déterminantes dans la définition des priorités du projet.

Je remercie également mon \textbf{encadrant académique, Prof.\ Adnane Elmrabty}, pour son encadrement scientifique, ses retours constructifs et son soutien lors de la rédaction de ce rapport. Sa rigueur méthodologique a été une boussole précieuse tout au long de ce parcours.

Mes remerciements s'adressent à l'ensemble des \textbf{enseignants de l'École Supérieure de Technologie de Guelmim}, dont la formation en Science de Données et Systèmes Intelligents m'a fourni les compétences directement mobilisées dans ce projet~: programmation Python, apprentissage automatique, bases de données et développement logiciel.

Enfin, je remercie chaleureusement ma \textbf{famille et mes proches} pour leur soutien moral indéfectible, leur patience et leurs encouragements constants tout au long de cette formation.

\vspace{1.5cm}
\begin{flushright}
  \textit{Abdelhafid Ibn Mhand}\\
  \textit{Guelmim, Juillet 2026}
\end{flushright}

% ── Dédicace ────────────────────────────────────────────────
\newpage
\chapter*{Dédicace}
\addcontentsline{toc}{chapter}{Dédicace}
\thispagestyle{empty}

\vspace*{3cm}
\begin{center}
\begin{minipage}{0.72\textwidth}
\setlength{\parskip}{10pt}
\centering\itshape\large

Je dédie ce modeste travail à mes \textbf{chers parents}, pour leur amour inconditionnel, leurs sacrifices silencieux et leur soutien permanent tout au long de mon parcours académique.

\vspace{0.5cm}
À ma \textbf{famille}, pour sa confiance et sa présence dans chaque moment décisif de ma vie.

\vspace{0.5cm}
À mes \textbf{enseignants}, qui m'ont transmis bien plus que des connaissances : une façon de penser, de questionner et de construire.

\vspace{0.5cm}
À mes \textbf{amis et collègues}, pour les échanges, les débats et les moments partagés qui rendent l'apprentissage vivant et l'effort moins solitaire.

\vspace{1.5cm}
{\normalsize\upshape --- Abdelhafid Ibn Mhand}
\end{minipage}
\end{center}

% ── Résumé ──────────────────────────────────────────────────
\newpage
\chapter*{Résumé}
\addcontentsline{toc}{chapter}{Résumé}

Ce projet de fin d'études porte sur la conception et la réalisation d'une plateforme de \textit{data intelligence} appliquée à l'analyse et à la prédiction de la pauvreté au Maroc. Il a été mené dans le cadre d'un stage au sein de la Direction Régionale du Haut-Commissariat au Plan de Guelmim, et s'appuie sur deux sources de données officielles téléchargées depuis le portail du HCP (\url{https://www.hcp.ma}) : l'\textbf{ENCDM} (Enquête Nationale sur la Consommation et les Dépenses des Ménages, 2019--2020, environ 16 000 ménages) et le \textbf{RGPH} (Recensement Général de la Population et de l'Habitat, 2014, plus de 730 000 ménages).

Le projet est articulé en quatre axes majeurs~: le \textbf{prétraitement des données} (nettoyage, encodage, standardisation et imputation adaptative via KNN et LightGBM), l'\textbf{analyse exploratoire} (concentration spatiale de la pauvreté, fracture urbain-rural, gradient éducatif, dimension de genre), la \textbf{modélisation prédictive} (LightGBM, hyperréseau PyTorch, apprentissage par transfert), et le \textbf{déploiement d'un tableau de bord Streamlit} interactif.

En parallèle, le projet a donné naissance à \textbf{Simpute}, une bibliothèque Python open source publiée sur \textbf{PyPI}, généralisant le moteur d'imputation adaptative développé pour la plateforme.

Les résultats clés révèlent que Guelmim-Oued Noun affiche le taux de pauvreté le plus élevé ($\approx$18\,\%), que l'écart de dépenses urbain-rural est de 2,3x, et que chaque palier d'éducation supplémentaire réduit le risque de pauvreté de 40 à 50\,\%. Le modèle LightGBM atteint un AUROC de 0,94 pour la prédiction de la pauvreté.

\vspace{0.4cm}
\noindent\textbf{Mots-clés :} Data Science, HCP, Pauvreté, Python, Streamlit, LightGBM, PyTorch, Imputation, Machine Learning, Simpute, PyPI, GitHub.

% ── Abstract ────────────────────────────────────────────────
\newpage
\chapter*{Abstract}
\addcontentsline{toc}{chapter}{Abstract}

This final year project focuses on the design and development of a data intelligence platform for poverty analysis and prediction in Morocco. It was carried out during an internship at the Regional Directorate of the High Commission for Planning (HCP) in Guelmim. Raw data were downloaded from the official HCP portal (\url{https://www.hcp.ma}) and encompass two institutional datasets: the \textbf{ENCDM} (National Household Consumption and Expenditure Survey, 2019--2020, approx.\ 16\,000 households) and the \textbf{RGPH} (General Population and Housing Census, 2014, over 730\,000 households).

The project is structured around four major phases: \textbf{data preprocessing} (cleaning, encoding, standardisation, and adaptive imputation using KNN and LightGBM); \textbf{exploratory analysis} (spatial poverty concentration, urban-rural divide, education gradient, gender dimension); \textbf{predictive modelling} (LightGBM baseline, PyTorch hypernetwork with geographic embeddings, transfer learning); and \textbf{Streamlit dashboard deployment} for interactive results exploration.

As a direct by-product of the internship, the project produced \textbf{Simpute}, an open-source Python library published on \textbf{PyPI}, which generalises the adaptive imputation engine built for the platform.

Key findings include: Guelmim-Oued Noun has the highest poverty rate ($\approx$18\,\%), the urban-rural expenditure gap is 2.3x, and each educational level reduces poverty risk by 40--50\,\%. The LightGBM classifier achieves an AUROC of 0.94.

\vspace{0.4cm}
\noindent\textbf{Keywords:} Data Science, HCP, Poverty, Python, Streamlit, LightGBM, PyTorch, Imputation, Machine Learning, Simpute, PyPI, GitHub, Cursor, Gemini.

% ── Table of Contents ───────────────────────────────────────
\newpage
\tableofcontents

\newpage
\listoffigures
\addcontentsline{toc}{chapter}{Liste des figures}

\newpage
\listoftables
\addcontentsline{toc}{chapter}{Liste des tableaux}

% ── Abbreviations ───────────────────────────────────────────
\newpage
\chapter*{Liste des abréviations}
\addcontentsline{toc}{chapter}{Liste des abréviations}

\begin{tabular}{@{}L{3.6cm}L{10.5cm}@{}}
\toprule
\textbf{Abréviation} & \textbf{Signification} \\
\midrule
HCP      & Haut-Commissariat au Plan \\
ENCDM    & Enquête Nationale sur la Consommation et les Dépenses des Ménages \\
RGPH     & Recensement Général de la Population et de l'Habitat \\
DBT      & Diplôme Bachelor Universitaire de Technologie \\
SDSI     & Science de Données et Systèmes Intelligents \\
EST      & École Supérieure de Technologie \\
MAD      & Dirham Marocain \\
DAP      & Dépense Annuelle par Personne \\
IPM      & Indice de Pauvreté Multidimensionnelle \\
KNN      & K-Nearest Neighbors (K Plus Proches Voisins) \\
LGBM     & LightGBM (Light Gradient Boosting Machine) \\
AUROC    & Area Under the Receiver Operating Characteristic Curve \\
SHAP     & SHapley Additive exPlanations \\
ML       & Machine Learning (Apprentissage Automatique) \\
DL       & Deep Learning (Apprentissage Profond) \\
API      & Application Programming Interface \\
CSV / Parquet & Formats de fichiers de données tabulaires \\
JSON     & JavaScript Object Notation \\
PyPI     & Python Package Index (registre de paquets Python open source) \\
CI/CD    & Continuous Integration / Continuous Deployment \\
SPSS     & Statistical Package for the Social Sciences (format de fichier .sav) \\
\bottomrule
\end{tabular}

% ============================================================
%  BODY
% ============================================================
\newpage
\hypersetup{pageanchor=true}
\pagenumbering{arabic}
\setcounter{page}{1}

% ============================================================
%  INTRODUCTION GÉNÉRALE
% ============================================================
\chapter*{Introduction générale}
\addcontentsline{toc}{chapter}{Introduction générale}

\epigraph{\itshape ``Les chiffres ne sont pas abstraits. Un écart de dépenses de 2,3x entre ménages urbains et ruraux n'est pas une statistique --- c'est la photographie d'une inégalité structurelle ancrée dans l'infrastructure, l'éducation et le marché du travail.''}{}

\section*{Contexte général}

La transformation numérique de la statistique publique ouvre de nouvelles perspectives pour l'analyse des politiques sociales. Les institutions nationales comme le Haut-Commissariat au Plan (HCP) produisent des volumes considérables de données officielles qui, correctement traitées, peuvent éclairer les décisions de politique publique avec une précision inédite. Le portail \href{https://www.hcp.ma}{\texttt{hcp.ma}} rend publiquement accessibles ces données~; pourtant, le passage de la donnée brute à l'information utile reste un défi technique et analytique majeur.

\section*{Motivation du projet}

Ce projet est né d'un constat simple~: les fichiers statistiques du HCP --- enquêtes ménages au format SPSS, tableaux Excel complexes --- contiennent une richesse d'information considérable qui reste largement sous-exploitée, notamment à l'échelle régionale. La Direction Régionale de Guelmim, dont le territoire affiche l'un des taux de pauvreté les plus élevés du Maroc, se trouve dans une position particulièrement propice pour bénéficier d'un outil de valorisation de ces données.

\section*{Objectif et périmètre}

L'objectif de ce projet est de concevoir une \textbf{plateforme de data intelligence} couvrant l'intégralité de la chaîne analytique~: ingestion des données brutes, pipeline de prétraitement robuste, analyse exploratoire structurée, modélisation prédictive à deux niveaux (gradient boosting et deep learning), et déploiement sous forme de tableau de bord interactif.

En marge du projet principal, une bibliothèque Python open source --- \textbf{Simpute} --- a été extraite du moteur d'imputation développé et publiée sur \textbf{PyPI}, constituant ainsi une contribution à la communauté Data Science au-delà du cadre du stage.

\section*{Outils et environnement de développement}

L'ensemble du projet a été développé en Python 3.10+, avec \textbf{Cursor} comme environnement de développement intégré (IDE) à assistance IA pour accélérer l'écriture et la révision du code. L'assistance de \textbf{Gemini} (Google) a été utilisée ponctuellement pour la recherche documentaire et la validation d'approches algorithmiques. Le versioning et la collaboration sont assurés via \textbf{GitHub}, avec des branches dédiées par fonctionnalité. La distribution open source de Simpute est gérée via \textbf{PyPI} et un workflow d'automatisation GitHub Actions.

\section*{Structure du rapport}

Ce document est organisé en neuf chapitres et deux annexes~:

\medskip
\begin{tabular}{@{}ll@{}}
\textbf{Chapitre 1} & Présentation de l'organisme d'accueil \\
\textbf{Chapitre 2} & Contexte du projet et analyse des besoins \\
\textbf{Chapitre 3} & Architecture générale et choix technologiques \\
\textbf{Chapitre 4} & Prétraitement des données \\
\textbf{Chapitre 5} & Analyse exploratoire et conclusions analytiques \\
\textbf{Chapitre 6} & Modélisation prédictive \\
\textbf{Chapitre 7} & Déploiement du tableau de bord Streamlit \\
\textbf{Chapitre 8} & Simpute~: de l'internship à l'open source \\
\textbf{Chapitre 9} & Résultats, apports et discussion \\
\textbf{Conclusion} & Synthèse et perspectives \\
\end{tabular}

% ============================================================
%  CHAPITRE 1 — ORGANISME D'ACCUEIL
% ============================================================
\chapter{Présentation de l'organisme d'accueil}

\section{Le Haut-Commissariat au Plan}

Le \textbf{Haut-Commissariat au Plan (HCP)} est l'institution nationale chargée de la production, de l'analyse et de la diffusion des statistiques officielles au Maroc. Placé directement sous la tutelle royale, il constitue la référence institutionnelle pour l'ensemble des données démographiques, économiques et sociales utilisées dans la planification et l'évaluation des politiques publiques du Royaume.

Son portail officiel (\href{https://www.hcp.ma}{\texttt{hcp.ma}}) est la source principale de données statistiques pour ce projet~: c'est depuis ce portail que les fichiers ENCDM et RGPH ont été téléchargés dans leurs versions brutes au format SPSS. Cette accessibilité publique des données constitue en elle-même un enjeu~: rendre ces données exploitables par un plus grand nombre d'acteurs est précisément la raison d'être de la plateforme développée.

Parmi les missions principales du HCP~:
\begin{itemize}
  \item Réalisation du \textbf{Recensement Général de la Population et de l'Habitat (RGPH)}, effectué tous les dix ans ;
  \item Production des \textbf{comptes nationaux} et calcul du Produit Intérieur Brut ;
  \item Élaboration d'indicateurs économiques et sociaux de référence nationale ;
  \item Conduite d'enquêtes auprès des ménages (ENCDM, HBS, etc.) et des entreprises ;
  \item Diffusion des statistiques officielles et appui aux administrations publiques.
\end{itemize}

\includegraphics[width=16\linewidth,height=10cm,keepaspectratio]{organigramme_hcp.png}

\section{Organisation et réseau territorial}

Le HCP est structuré autour d'une \textbf{administration centrale} basée à Rabat et d'un réseau de \textbf{directions régionales} couvrant les douze régions du Maroc. L'administration centrale définit les orientations stratégiques, coordonne les opérations statistiques nationales et produit les indicateurs agrégés. Les directions régionales, quant à elles, assurent la collecte sur le terrain, le suivi des enquêtes locales et la diffusion des statistiques à l'échelle territoriale.

Cette décentralisation est fondamentale~: elle garantit que les particularités de chaque territoire --- dynamiques démographiques, contraintes économiques, besoins spécifiques --- sont correctement captées et rapportées.
Dans le cas de Guelmim-Oued Noun, région à la fois étendue géographiquement et fragile économiquement, ce maillage territorial revêt une importance particulière.

\section{La Direction Régionale de Guelmim-Oued Noun}

La \textbf{Direction Régionale du HCP de Guelmim-Oued Noun} est l'entité qui a accueilli ce stage. Son périmètre d'intervention couvre les quatre provinces de la région~: \textbf{Guelmim}, \textbf{Assa-Zag}, \textbf{Tan-Tan} et \textbf{Sidi Ifni}.

\begin{table}[H]
\centering
\caption{Activités et missions de la Direction Régionale du HCP de Guelmim}
\label{tab:hcp_services}
\begin{tabularx}{\textwidth}{L{3.8cm}X}
\toprule
\textbf{Activité} & \textbf{Description} \\
\midrule
Collecte de données & Coordination des opérations de collecte~: enquêtes ménages, RGPH, recensements économiques \\
Traitement statistique & Validation, codification et traitement des données collectées sur le terrain \\
Production d'indicateurs & Élaboration d'indicateurs régionaux sur la démographie, l'économie, le social et le logement \\
Publication et diffusion & Rapports statistiques régionaux, notes de synthèse, tableaux de données accessibles en ligne \\
Appui institutionnel & Assistance technique aux administrations régionales, collectivités territoriales et chercheurs \\
\bottomrule
\end{tabularx}
\end{table}

\section{Contexte socio-économique de la région}

La région Guelmim-Oued Noun occupe une position stratégique dans le territoire national, à la jonction du Maroc central et des provinces du Grand Sud. Ses caractéristiques socio-économiques en font un terrain d'étude particulièrement pertinent~:

\begin{itemize}
  \item \textbf{Taux de pauvreté parmi les plus élevés du Royaume}~: les données ENCDM indiquent un taux d'environ 18\,\%, significativement au-dessus de la moyenne nationale ;
  \item \textbf{Fracture urbain-rural marquée}~: les conditions de vie, l'accès aux services et les revenus des ménages diffèrent substantiellement entre les centres urbains comme Guelmim et les zones rurales des provinces d'Assa-Zag et de Sidi Ifni ;
  \item \textbf{Faible densité d'infrastructure numérique}~: l'accès à internet des ménages ruraux de la région est estimé à moins de 15\,\%, contre 48\,\% en milieu urbain au niveau national ;
  \item \textbf{Économie largement informelle}~: une part significative de la population active opère en dehors des circuits formels d'emploi, rendant les données d'enquête d'autant plus précieuses pour cerner la réalité économique.
\end{itemize}

C'est dans ce contexte que la Direction Régionale du HCP a souhaité disposer d'un outil permettant de transformer les données statistiques disponibles en informations directement exploitables pour la compréhension du territoire et l'aide à la décision.

\section{Contexte et déroulement du stage}

Le stage s'est déroulé au sein de la Direction Régionale sur une durée de plusieurs semaines. La mission assignée consistait à concevoir une plateforme analytique permettant de valoriser les données officielles disponibles, avec un focus sur l'analyse de la pauvreté. Le stage a également été l'occasion de collaborer avec les équipes de la direction pour comprendre les modes d'utilisation actuels des données et identifier les besoins réels en matière de visualisation et d'aide à la décision.

L'environnement de travail a combiné les locaux de la direction régionale pour la compréhension du domaine et les échanges avec l'encadrant professionnel, et un environnement de développement autonome --- équipé de \textbf{Cursor}, \textbf{Python 3.10+} et \textbf{Git/GitHub} --- pour la réalisation technique du projet.

% ============================================================
%  CHAPITRE 2 — CONTEXTE ET BESOINS
% ============================================================
\chapter{Contexte du projet et analyse des besoins}

\section{Problématique générale}

La disponibilité des données statistiques officielles ne suffit pas à produire de la connaissance opérationnelle. Le HCP génère des volumes considérables d'informations~: enquêtes ménages, recensements, comptes régionaux. Mais ces données sont souvent dispersées dans des formats hétérogènes (SPSS, Excel), présentées sous forme de tableaux complexes à structure hiérarchique, et rarement mises en relation les unes avec les autres pour révéler les dynamiques structurelles sous-jacentes.

\begin{infobox}[Problématique centrale du projet]
Comment exploiter conjointement les données de l'ENCDM et du RGPH --- téléchargées depuis \texttt{hcp.ma} --- pour construire une plateforme analytique permettant (1) de comprendre les déterminants structurels de la pauvreté au Maroc, (2) de prédire le statut de pauvreté des ménages avec haute précision, et (3) de diffuser ces résultats via un tableau de bord interactif accessible aux décideurs de la Direction Régionale~?
\end{infobox}

\section{Sources de données officielles}

\subsection{Acquisition via le portail HCP}

Les deux jeux de données utilisés dans ce projet ont été téléchargés depuis le portail officiel du Haut-Commissariat au Plan (\href{https://www.hcp.ma}{\texttt{hcp.ma}}). Ce portail constitue la source de référence pour les statistiques nationales et régionales du Maroc, et met à disposition des chercheurs et des institutions des microdonnées anonymisées sous forme de fichiers SPSS.

\subsection{ENCDM 2019--2020}

L'\textbf{Enquête Nationale sur la Consommation et les Dépenses des Ménages} est l'instrument principal utilisé par le HCP pour mesurer la pauvreté monétaire. Elle porte sur environ \textbf{16~000 ménages} représentatifs de la population nationale, avec des coefficients de pondération (poids de sondage) permettant de généraliser les estimations.

Les variables clés incluent~: dépenses annuelles par ménage et par personne (DAP), quintiles et déciles de dépense, statut de pauvreté et de vulnérabilité, milieu de résidence (urbain/rural), région, niveau d'éducation du chef de ménage, profession et secteur d'activité, taille et structure démographique du ménage.

\subsection{RGPH 2014}

Le \textbf{Recensement Général de la Population et de l'Habitat} est un recensement exhaustif couvrant plus de \textbf{730~000 ménages}. Il ne contient pas de données de dépenses ni de statut de pauvreté, mais fournit une couverture quasi-complète des conditions d'habitat et d'infrastructure, indispensable pour l'apprentissage par transfert et l'enrichissement contextuel des modèles.

Les variables clés décrivent~: type de logement et matériaux de construction, accès à l'eau potable, à l'électricité et à internet, possession de téléphones mobiles et de véhicules, composition et taille du ménage.

\section{Objectifs du projet}

\begin{enumerate}
  \item \textbf{Prétraitement robuste}~: Nettoyer, encoder, standardiser et imputer les données ENCDM et RGPH de manière reproductible et documentée.
  \item \textbf{Analyse exploratoire structurée}~: Identifier les déterminants structurels de la pauvreté (géographie, éducation, genre, infrastructure).
  \item \textbf{Modélisation prédictive}~: Construire des classifieurs de haute précision pour prédire le statut de pauvreté des ménages.
  \item \textbf{Déploiement interactif}~: Développer un tableau de bord Streamlit permettant l'exploration des résultats et la prédiction en temps réel.
  \item \textbf{Contribution open source}~: Extraire et publier le moteur d'imputation sous forme de bibliothèque Python réutilisable sur PyPI.
\end{enumerate}

\section{Besoins fonctionnels}

\begin{table}[H]
\centering
\caption{Besoins fonctionnels de la plateforme}
\label{tab:besoins_fonc}
\begin{tabularx}{\textwidth}{C{1.2cm}L{4.2cm}X}
\toprule
\textbf{ID} & \textbf{Besoin} & \textbf{Description détaillée} \\
\midrule
BF01 & Ingestion des données & Chargement des fichiers SPSS (.sav) et transformation en DataFrames pandas manipulables \\
BF02 & Prétraitement automatique & Encodage catégoriel, standardisation des features continues et imputation des valeurs manquantes \\
BF03 & Analyse exploratoire & Visualisation des distributions, taux de pauvreté par région, éducation, genre et infrastructure \\
BF04 & Prédiction interactive & Saisie des caractéristiques d'un ménage et obtention d'une probabilité de pauvreté en temps réel \\
BF05 & Comparaison de modèles & Affichage simultané des résultats LightGBM et du hyperréseau pour comparer les architectures \\
BF06 & Exploration des données & Interface paginée et filtrable pour consulter les données brutes et imputées \\
BF07 & Reproductibilité & Persistance des scalers, imputers et mappings pour garantir des résultats identiques à chaque exécution \\
\bottomrule
\end{tabularx}
\end{table}

\section{Besoins non fonctionnels}

\begin{itemize}
  \item \textbf{Performance}~: Chargement initial du dashboard en moins de 10 secondes~; navigation entre onglets en moins d'une seconde grâce au caching Streamlit.
  \item \textbf{Reproductibilité}~: Pipeline idempotent -- toute ré-exécution produit des résultats identiques grâce aux artefacts persistés sur disque.
  \item \textbf{Maintenabilité}~: Architecture modulaire avec séparation stricte des responsabilités entre les couches données, analyse, modèles et interface.
  \item \textbf{Portabilité}~: Application déployable localement sans serveur externe ; modèles conçus pour fonctionner exclusivement en CPU.
  \item \textbf{Traçabilité}~: Chaque décision analytique documentée dans les notebooks Jupyter avec justification explicite.
\end{itemize}

\section{Méthodologie adoptée}

La méthodologie suit un cycle itératif inspiré des meilleures pratiques en Data Science industrielle, structuré en quatre phases progressives~:

\includegraphics[width=5\linewidth,height=5.5cm,keepaspectratio]{methodologie_projet.png}

\textbf{Phase 1 -- Compréhension des données.} Lecture critique des fichiers ENCDM et RGPH, identification des variables pertinentes, cartographie des valeurs manquantes et formulation des hypothèses analytiques.

\textbf{Phase 2 -- Prétraitement.} Nettoyage, encodage, standardisation et imputation. Cette phase est la plus longue et la plus critique~: des données mal structurées contaminent irrémédiablement toutes les étapes suivantes.

\textbf{Phase 3 -- Analyse et modélisation.} Analyse exploratoire pour dégager les patterns structurels, puis construction et évaluation des modèles prédictifs.

\textbf{Phase 4 -- Déploiement.} Intégration de l'ensemble dans un tableau de bord Streamlit et publication de Simpute sur PyPI.

% ============================================================
%  CHAPITRE 3 — ARCHITECTURE ET TECHNOLOGIES
% ============================================================
\chapter{Architecture générale et choix technologiques}

\section{Vue d'ensemble de l'architecture}

La plateforme est organisée selon une \textbf{architecture en couches} stricte, où chaque couche est responsable d'un niveau d'abstraction précis. Cette séparation rend le projet maintenable, testable et évolutif.

\includegraphics[width=5\linewidth,height=5cm,keepaspectratio]{general.png}

Les quatre couches principales sont~:
\begin{enumerate}
  \item \textbf{Couche données}~: fichiers SPSS bruts téléchargés depuis \texttt{hcp.ma}, transformés en fichiers Parquet nettoyés.
  \item \textbf{Couche traitement}~: notebooks Jupyter orchestrant le pipeline de prétraitement, d'analyse et de modélisation.
  \item \textbf{Couche artefacts}~: scalers, imputers, mappings et checkpoints de modèles persistés sur disque en formats joblib, JSON et PyTorch.
  \item \textbf{Couche présentation}~: modules Python du dashboard Streamlit, chacun responsable d'une fonctionnalité spécifique.
\end{enumerate}

\section{Environnement et outils de développement}

\begin{table}[H]
\centering
\caption{Environnement et outils de développement}
\label{tab:devtools}
\begin{tabularx}{\textwidth}{L{3.5cm}L{3cm}X}
\toprule
\textbf{Outil} & \textbf{Catégorie} & \textbf{Rôle dans le projet} \\
\midrule
\textbf{Cursor} & IDE / Éditeur IA & Environnement de développement principal avec assistance IA intégrée pour l'écriture, la révision et le débogage du code \\
\textbf{Gemini} & Assistant IA & Recherche documentaire, validation d'approches algorithmiques et génération de documentation \\
\textbf{GitHub} & Versioning & Hébergement du code, branches dédiées par fonctionnalité, workflow CI/CD pour la publication PyPI \\
\textbf{PyPI} & Distribution & Publication et distribution open source de la bibliothèque Simpute à la communauté Python \\
\textbf{hcp.ma} & Source de données & Portail officiel du HCP pour le téléchargement des microdonnées ENCDM et RGPH \\
\textbf{Jupyter} & Exploration & Notebooks pour le prototypage, l'analyse interactive et la documentation des décisions analytiques \\
\bottomrule
\end{tabularx}
\end{table}

\section{Pile technologique Python}

\begin{table}[H]
\centering
\small
\caption{Pile technologique Python du projet}
\label{tab:techstack}
\begin{tabularx}{\textwidth}{@{}L{2.8cm} L{2.8cm} >{\raggedright\arraybackslash}X@{}}
\toprule
\textbf{Composant} & \textbf{Technologie} & \textbf{Justification} \\
\midrule
Langage principal & Python 3.10+ & Écosystème Data Science le plus mature, lecture native SPSS via pyreadstat \\
Manipulation & Pandas, NumPy & DataFrames haute performance, opérations vectorisées sur grands volumes \\
Imputation & scikit-learn, LightGBM & KNN pour petits jeux de données (ENCDM), LGBM pour grands volumes (RGPH) \\
Modélisation & LightGBM, PyTorch & LGBM état de l'art sur données tabulaires, PyTorch pour le hyperréseau expérimental \\
Visualisation & Plotly, Matplotlib & Graphiques interactifs (Plotly) pour le dashboard, statiques (Matplotlib) pour les rapports \\
Dashboard & Streamlit & Déploiement rapide, composants réactifs, caching natif \\
Persistance & Parquet, Joblib, JSON & Lecture 10--100x plus rapide que CSV, préservation des types, sérialisation ML \\
Versioning & Git, GitHub & Historique complet, branches dédiées, intégration CI/CD \\
Distribution & PyPI, pyproject.toml & Publication standardisée de Simpute, workflow de release automatisé \\
\bottomrule
\end{tabularx}
\end{table}

\section{Choix de Cursor comme IDE}

\textbf{Cursor} est un éditeur de code à assistance IA (basé sur VS Code) qui a joué un rôle central dans la productivité du projet. Contrairement à un IDE classique, Cursor permet de~:

\begin{itemize}
  \item Générer des ébauches de code complexe (pipelines d'imputation, architecture du hyperréseau) à partir de descriptions en langage naturel ;
  \item Refactoriser automatiquement du code répétitif, comme les boucles d'imputation colonne par colonne du RGPH ;
  \item Détecter proactivement les erreurs de logique, notamment dans la gestion des graphes de dépendance lors de l'imputation KNN ;
  \item Générer des docstrings et de la documentation de manière cohérente à travers tous les modules.
\end{itemize}

L'utilisation de Cursor n'a pas remplacé la réflexion algorithmique -- les choix d'architecture, les décisions de modélisation et les interprétations analytiques restent entièrement celles du développeur -- mais elle a significativement réduit le temps passé sur les tâches de codage répétitives, permettant de concentrer l'effort sur la résolution de problèmes réellement complexes.

\section{Rôle de Gemini dans la recherche documentaire}

\textbf{Gemini} (Google DeepMind) a été mobilisé de manière ponctuelle, essentiellement pour~:

\begin{itemize}
  \item La recherche bibliographique sur les techniques d'imputation avancées et les architectures de hyperréseaux ;
  \item La validation de choix algorithmiques~: par exemple, comparer les stratégies d'imputation KNN vs MICE vs gradient boosting pour des données de recensement ;
  \item La génération de formulations alternatives pour les sections analytiques les plus techniques du rapport.
\end{itemize}

\section{Organisation du dépôt GitHub}

\begin{lstlisting}[style=treenostyle, caption={Structure du depot HCP-Internship (branche Dashboard)}]
HCP-Internship/
|-- Notebooks/
|   |-- Pre-processing.ipynb
|   |-- Analysis.ipynb
|   |-- Modeling.ipynb
|-- Dashboard/
|   |-- app.py
|   |-- data_loader.py
|   |-- plots.py
|   |-- mapping.py
|   |-- sandbox.py
|   |-- hypernet.py
|   |-- theme.py
|   |-- components.py
|-- Artifacts/
|   |-- scalers/
|   |-- imputers/
|   |-- mappings/
|   |-- models/
|-- Data/
|   |-- CleanENCDM.parquet
|   |-- CleanRGPH.parquet
|-- requirements.txt
|-- README.md
\end{lstlisting}

% ============================================================
%  CHAPITRE 4 — PRÉTRAITEMENT
% ============================================================
\chapter{Prétraitement des données}

\section{Vue d'ensemble et principes directeurs}

Le prétraitement des données est l'étape la plus critique et la plus chronophage du projet. Son importance est souvent sous-estimée~: des données mal structurées ou incohérentes produisent invariablement des analyses erronées, quelle que soit la sophistication des modèles appliqués en aval. Cette réalité est particulièrement prononcée avec les données d'enquêtes publiques, qui sont conçues pour la consultation humaine et non pour le traitement automatique.

Le pipeline a été conçu selon trois principes directeurs~:
\begin{itemize}
  \item \textbf{Idempotence}~: toute ré-exécution complète produit un résultat identique au bit près.
  \item \textbf{Auditabilité}~: chaque décision de transformation est documentée dans le notebook avec sa justification.
  \item \textbf{Séparabilité}~: les artefacts (scalers, imputers, mappings) sont persistés séparément du code, pour réutilisation dans le dashboard sans ré-exécution.
\end{itemize}

\includegraphics[width=5\linewidth,height=5.5cm,keepaspectratio]{process.png}

\section{Ingestion et inspection initiale}

La première étape consiste à charger les fichiers SPSS depuis le disque via la bibliothèque \texttt{pyreadstat}. Celle-ci préserve les métadonnées du fichier SPSS (labels de variables, codes de valeur) qui seront utiles pour la construction des mappings de décodage.

\begin{table}[H]
\centering
\small
\caption{Métriques initiales des jeux de données brutes}
\label{tab:initial_metrics}
\begin{tabularx}{\textwidth}{@{}lrrrX@{}}
\toprule
\textbf{Jeu de données} & \textbf{Lignes} & \textbf{Colonnes} & \textbf{Nulls} & \textbf{Format} \\
\midrule
ENCDM & $\approx$16\,000 & 26 & $\approx$3.2\,\% & SPSS (.sav) \\
RGPH  & $\approx$730\,000 & 36 & Bimodal (0--15\,\%) & SPSS (.sav) \\
\bottomrule
\end{tabularx}
\end{table}

Le profil de missingness du RGPH est bimodal~: certaines colonnes sont quasi-complètes (accès à l'électricité, taille du ménage), d'autres présentent des taux élevés (possession d'un véhicule pour les ménages sans véhicule). Ce pattern caractéristique des données de recensement impose des stratégies d'imputation distinctes par colonne --- c'est précisément ce qui a motivé le développement de Simpute.

\section{Résolution d'alias et harmonisation des noms}

Les données brutes du HCP utilisent des noms de variables en français avec des accents incohérents et des variantes orthographiques multiples selon les millésimes d'enquête. Un \textbf{résolveur d'alias} centralise ces variations~:

\begin{lstlisting}[caption={Resolution d'alias -- normalisation des noms de colonnes}]
ALIAS_MAP = dict(
    Region_12=["Region_12", "region_12", "REGION_12"],
    Vulnerable=["Vulnerable", "vulnerable", "VULNERABLE"],
    Pauvre=["Pauvre", "pauvre", "PAUVRE"],
)

def resolve_alias(col, alias_map):
    for canonical, variants in alias_map.items():
        if col in variants:
            return canonical
    return col

df.columns = [resolve_alias(c, ALIAS_MAP) for c in df.columns]
\end{lstlisting}

\section{Encodage catégoriel et registre des types}

Les colonnes catégorielles sont transformées en entiers via \texttt{pd.factorize}. Chaque mapping entier $\leftrightarrow$ étiquette est persisté en JSON sous \texttt{Artifacts/mappings/}. Ce registre garantit que le dashboard et les notebooks utilisent les mêmes décodages.

\begin{lstlisting}[caption={Encodage categoriel et persistance des mappings}]
import json, os
os.makedirs("Artifacts/mappings", exist_ok=True)

for col in categorical_cols:
    df[col], uniques = pd.factorize(df[col])
    mapping = dict(enumerate(uniques))
    path = "Artifacts/mappings/" + col + ".json"
    with open(path, "w") as f:
        json.dump(mapping, f, ensure_ascii=False, indent=2)
\end{lstlisting}

\section{Standardisation des features continues}

Les features continues sont standardisées individuellement via des objets \texttt{StandardScaler} dédiés~:

\begin{itemize}
  \item \textbf{Nécessité pour KNN}~: sans normalisation, les variables à grande amplitude (ex.\ DAP en MAD) domineraient la recherche de voisins.
  \item \textbf{Nécessité pour les réseaux de neurones}~: les features doivent être dans un ordre de grandeur comparable pour une convergence stable.
\end{itemize}

Chaque scaler est persisté en \texttt{.joblib} sous le nom de la colonne correspondante. Cela constitue un \textbf{contrat de transformation}~: lorsque le bac à sable prédictif du dashboard reçoit la saisie d'un utilisateur, il applique exactement le même scaler que celui utilisé lors de l'entraînement.

\section{Imputation des valeurs manquantes}

\subsection{Pourquoi imputer plutôt que supprimer~?}

La suppression des lignes incomplètes introduirait un \textbf{biais de sélection systématique}~: les ménages avec des valeurs manquantes ne sont pas aléatoirement distribués dans la population~; ils présentent des caractéristiques structurelles spécifiques (ruralité, taille, secteur d'activité informel) qui sont précisément les déterminants de la pauvreté que l'on cherche à analyser. Supprimer ces observations biaiserait donc à la fois l'analyse exploratoire et les modèles prédictifs.

\subsection{Imputation ENCDM --- K-Nearest Neighbors}

Pour l'ENCDM ($\approx$16\,000 lignes, features hétérogènes richement informatives), l'imputation KNN est appropriée~:

\begin{itemize}
  \item Elle exploite la similarité entre ménages pour ``emprunter'' les valeurs manquantes de voisins proches dans l'espace des features complètes.
  \item Des \textbf{graphes de dépendance} encodent la connaissance du domaine et préviennent la fuite d'information entre features non liées.
  \item La quantité de données (16\,000 lignes) est suffisante pour que les $k$ voisins trouvés soient représentatifs de la distribution réelle.
\end{itemize}

\subsection{Imputation RGPH --- LightGBM}

Pour le RGPH (730\,000 lignes, variables catégorielles à haute cardinalité), LightGBM est le choix optimal~:

\begin{itemize}
  \item Les ensembles d'arbres capturent les interactions non linéaires complexes que KNN manquerait dans un espace de 36 dimensions.
  \item LightGBM gère nativement les variables catégorielles, évitant l'encodage one-hot et ses problèmes de dimensionnalité.
  \item Les hyperparamètres sont ajustés \textbf{dynamiquement} selon le taux de missingness~: plus une colonne est lacunaire, plus le modèle reçoit d'estimateurs (pour mieux généraliser) et des feuilles plus profondes (pour capturer plus d'interactions).
\end{itemize}

\begin{lstlisting}[caption={Pipeline d'imputation LightGBM pour le RGPH -- version simplifiee}]
import joblib
from lightgbm import LGBMClassifier

for target_col in rgph_missing_cols:
    artifact_path = "Artifacts/imputers/RGPH_" + target_col + ".joblib"

    if os.path.exists(artifact_path):
        imputer = joblib.load(artifact_path)
    else:
        X_train = df.dropna(subset=[target_col]).drop(columns=[target_col])
        y_train = df.dropna(subset=[target_col])[target_col]

        n_missing = df[target_col].isna().sum()
        imputer = LGBMClassifier(
            n_estimators=max(100, int(n_missing / 100)),
            num_leaves=int(pow(n_missing, 0.4)),
            min_child_samples=5,
        ).fit(X_train, y_train)
        joblib.dump(imputer, artifact_path)

    mask = df[target_col].isna()
    df.loc[mask, target_col] = imputer.predict(
        df[mask].drop(columns=[target_col]))
\end{lstlisting}

\section{Export et validation}

Après imputation, les DataFrames nettoyés sont exportés au format \textbf{Parquet}. Ce format a été préféré au CSV pour trois raisons concrètes~: (1) lecture 10 à 100 fois plus rapide, (2) préservation native des types de données sans re-parsing, (3) support de la sélection columnar -- le dashboard peut charger uniquement les colonnes dont il a besoin.

\begin{table}[H]
\centering
\caption{Qualité des données après le pipeline complet}
\label{tab:data_quality}
\begin{tabular}{lrrrr}
\toprule
\textbf{Dataset} & \textbf{Observations} & \textbf{Features finales} & \textbf{Taux de nulls} & \textbf{Format} \\
\midrule
CleanENCDM & 15\,970 & 78 & $\approx$2.3\,\% & Parquet \\
CleanRGPH  & 730\,099 & 36 & $\approx$1.1\,\% & Parquet \\
\bottomrule
\end{tabular}
\end{table}

Une \textbf{validation aller-retour} est effectuée~: chaque fichier Parquet est rechargé depuis le disque et comparé au DataFrame source pour détecter toute corruption ou perte de précision lors de la sérialisation. Cette étape, codée en moins de dix lignes, a permis de détecter en phase de développement une incompatibilité de type entre pandas et Arrow pour certaines colonnes entières.

% ============================================================
%  CHAPITRE 5 — ANALYSE EXPLORATOIRE
% ============================================================
\chapter{Analyse exploratoire et conclusions analytiques}

\section{Approche et précautions méthodologiques}

L'analyse exploratoire a été conduite dans le notebook \texttt{Analysis.ipynb}. Les features standardisées ont été \textbf{inversement transformées} vers leurs unités originales. Les pondérations de sondage (\texttt{coef\_indiv}, \texttt{coef\_menage}) ont été appliquées systématiquement.

\section{Concentration spatiale de la pauvreté}

Le premier résultat majeur est que la pauvreté n'est pas distribuée uniformément sur le territoire marocain. Elle est fortement concentrée dans les régions du sud et de l'est, tandis que les régions atlantiques et du nord-ouest affichent des taux nettement inférieurs.

\begin{table}[H]
\centering
\small
\caption{Taux de pauvreté par région -- ENCDM 2019-2020 (estimations pondérées)}
\label{tab:poverty_region}
\begin{tabularx}{\textwidth}{@{}l r X@{}}
\toprule
\textbf{Région} & \textbf{Taux} & \textbf{Caractérisation} \\
\midrule
Guelmim-Oued Noun & $\approx$18\,\% & Le plus élevé -- région hôte du stage \\
Draa-Tafilalet & $\approx$15\,\% & Désertique, infrastructure limitée \\
Béni Mellal-Khénifra & $\approx$14\,\% & Enclavement rural persistant \\
Rabat-Salé-Kénitra & $\approx$8\,\% & Sous la moyenne -- capitale administrative \\
Casablanca-Settat & $\approx$6.5\,\% & Le plus bas -- capital économique \\
\bottomrule
\end{tabularx}
\end{table}

\includegraphics[width=10\linewidth,height=5.cm,keepaspectratio]{carte.png}
\includegraphics[width=10\linewidth,height=5.5cm,keepaspectratio]{newplot.png}

\section{La fracture urbain-rural}

Si l'on devait identifier le \textbf{pattern statistique dominant} des données, ce serait sans conteste la fracture urbain-rural. Les chiffres sont frappants~:

\begin{resultbox}
\textbf{Dépenses annuelles moyennes par ménage (ENCDM, pondéré)~:}\\[4pt]
\hspace{1em}$\bullet$\quad Ménages urbains : 45\,000 MAD/an\\
\hspace{1em}$\bullet$\quad Ménages ruraux : 22\,000 MAD/an\\
\hspace{1em}$\bullet$\quad \textbf{Rapport : 2,3x}\\[6pt]
\textbf{Taux de pauvreté par milieu~:}\\[4pt]
\hspace{1em}$\bullet$\quad Milieu urbain : 8.2\,\%\\
\hspace{1em}$\bullet$\quad Milieu rural : 19.4\,\%\\
\hspace{1em}$\bullet$\quad \textbf{Rapport : 2.4x}
\end{resultbox}

Ces deux ratios -- dépenses et pauvreté -- progressent en parfaite synchronie, confirmant que la ruralité est le prédicteur structurel dominant de la vulnérabilité économique. Ce n'est pas seulement une question de géographie per se, mais de ce que la géographie détermine~: accès à l'électricité (99\,\% urbain vs 94\,\% rural selon le RGPH), à internet (48\,\% vs 18\,\%), aux marchés formels d'emploi, et à l'offre éducative secondaire et supérieure.

\includegraphics[width=14\linewidth,height=8.5cm,keepaspectratio]{urur.png}

\section{L'éducation comme levier de réduction de la pauvreté}

L'analyse croisée du niveau d'éducation et du statut de pauvreté révèle un gradient remarquablement régulier et d'une ampleur peu commune~:

\begin{table}[H]
\centering
\caption{Taux de pauvreté selon le niveau d'éducation du chef de ménage}
\label{tab:educ_poverty}
\begin{tabular}{lrl}
\toprule
\textbf{Niveau d'éducation} & \textbf{Taux de pauvreté} & \textbf{Réduction vs niveau précédent} \\
\midrule
Sans scolarité & 24.1\,\% & --- (référence) \\
Niveau primaire & 16.8\,\% & $-$30\,\% \\
Niveau secondaire & 9.2\,\% & $-$45\,\% \\
Enseignement supérieur & 3.4\,\% & $-$63\,\% \\
\bottomrule
\end{tabular}
\end{table}

Chaque palier d'éducation réduit le risque de pauvreté de 40 à 50\,\%. Ce n'est pas une simple corrélation~: l'éducation opère via plusieurs canaux simultanés -- accès à l'emploi formel, salaires plus élevés, meilleure information sur les opportunités, réseaux sociaux plus étendus et capacité à naviguer les systèmes d'aide sociale. L'implication politique directe est que l'expansion de l'accès à l'enseignement secondaire dans les zones rurales constituerait la stratégie de réduction de la pauvreté à rendement potentiel le plus élevé.

\includegraphics[width=16\linewidth,height=10cm,keepaspectratio]{education.png}

\section{La dimension de genre}

Les ménages dirigés par une femme affichent un taux de pauvreté de \textbf{14.8\,\%} contre \textbf{11.2\,\%} pour les ménages dirigés par un homme, soit un écart de 3.6 points de pourcentage. Cet écart persiste après contrôle du niveau d'éducation, du secteur d'activité et de la région, ce qui indique que le genre opère comme une barrière structurelle indépendante.

Les mécanismes sous-jacents sont multiples~: discrimination salariale dans l'emploi formel, accès plus limité au crédit et à la propriété foncière, responsabilités de soin inégalement réparties qui contraignent la participation au marché du travail, et réseaux professionnels plus étroits. Ces facteurs se combinent pour produire un désavantage économique systémique qui ne se résorbe pas spontanément avec l'élévation du niveau d'éducation.

\section{Qualité du logement comme proxy du niveau de vie}

À partir des données du RGPH, un \textbf{indice composite de qualité du logement} a été construit en agrégeant plusieurs dimensions~: type d'habitat (moderne vs traditionnel), matériaux des murs, accès à l'eau courante, à l'électricité et à internet, possession d'appareils électroménagers. Cet indice corrèle à \textbf{r = 0.71} avec les taux de pauvreté régionaux calculés depuis l'ENCDM.

Cette corrélation élevée entre deux sources de données indépendantes (RGPH et ENCDM) valide la cohérence des deux enquêtes et ouvre une perspective opérationnelle importante~: dans les contextes où les données de dépenses sont indisponibles, les indicateurs de logement du recensement peuvent servir de proxy robuste pour estimer les niveaux de pauvreté régionaux.

\includegraphics[width=5\linewidth,height=5cm,keepaspectratio]{longement.png}

\section{Synthèse analytique}

Les quatre déterminants identifiés -- géographie, éducation, genre, infrastructure -- ne sont pas indépendants~; ils forment un système d'inégalités imbriquées~: la ruralité limite l'accès à l'éducation, qui contraint l'emploi formel, qui reproduit la pauvreté intergénérationnellement. Les ménages ruraux dirigés par des femmes sans scolarité cumulent les vulnérabilités de chaque dimension. C'est cette nature systémique qui justifie une approche multivariée de la prédiction plutôt que des indicateurs univariés.

% ============================================================
%  CHAPITRE 6 — MODÉLISATION
% ============================================================
\chapter{Modélisation prédictive}

\section{Architecture duale et philosophie de la modélisation}

La démarche de modélisation est structurée autour de trois questions hiérarchisées~:
\begin{enumerate}
  \item Peut-on prédire la pauvreté avec haute précision à l'aide de méthodes d'apprentissage automatique standard~?
  \item Une architecture de deep learning intégrant le contexte géographique améliore-t-elle les performances~?
  \item Quelles features pilotent les prédictions, et sont-elles cohérentes avec les intuitions analytiques du chapitre précédent~?
\end{enumerate}

Ces questions ont conduit à la mise en place de deux moteurs d'inférence parallèles, conçus pour répondre conjointement et offrir une comparaison rigoureuse.

\section{Moteur 1 : LightGBM Baseline}

LightGBM a été choisi comme modèle de référence pour quatre raisons concrètes~:

\begin{itemize}
  \item \textbf{Dominance sur les données tabulaires}~: les ensembles d'arbres à gradient boosté constituent l'état de l'art pour les données tabulaires hétérogènes.
  \item \textbf{Support natif des poids de sondage}~: LightGBM accepte directement les coefficients \texttt{coef\_menage}, essentiels pour une inférence non biaisée sur l'ENCDM.
  \item \textbf{Interprétabilité via SHAP}~: les importances de features et les valeurs SHAP permettent d'expliquer chaque prédiction individuelle, une exigence implicite dans un contexte institutionnel.
  \item \textbf{Performance d'inférence}~: inférence en millisecondes, compatible avec un usage interactif dans le dashboard.
\end{itemize}

Des classifieurs séparés ont été entraînés pour la pauvreté (\texttt{Pauvre}) et la vulnérabilité (\texttt{Vulnerable}), optimisés sur le score F1 pour gérer le déséquilibre naturel des classes (les ménages pauvres sont minoritaires dans l'échantillon).

\section{Moteur 2 : Hyperréseau PyTorch}

L'\textbf{hyperréseau} est la composante expérimentale du projet. L'hypothèse centrale est la suivante~: un classifieur LightGBM standard traite toutes les régions de manière identique. Si l'on apprend des embeddings denses pour la région et le milieu, et que l'on utilise ces embeddings pour \textit{générer les poids} d'un prédicteur aval, le modèle pourrait adapter ses frontières de décision au contexte géographique local.

\includegraphics[width=16\linewidth,height=10cm,keepaspectratio]{hyper.png}

\textbf{Composants de l'architecture~:}
\begin{itemize}
  \item \texttt{RgphEmbedding}~: projette les codes catégoriels du RGPH (région, province, milieu, type de logement) dans des vecteurs denses de dimension 16.
  \item \texttt{EncdmEmbedding}~: projette les codes catégoriels de l'ENCDM dans des vecteurs de dimension 8.
  \item \texttt{Hypernet}~: réseau générateur qui prend la concaténation des embeddings et produit la matrice de poids du prédicteur final, lui permettant de s'adapter au contexte géographique.
\end{itemize}

Le modèle a été entraîné exclusivement en CPU pour garantir la portabilité. Le checkpoint \texttt{Hypernet.pt} stocke les poids, les dimensions d'embedding, le seuil de décision calibré et l'historique complet de la loss par époque.

\section{Résultats comparatifs}

\begin{table}[H]
\centering
\caption{Comparaison des performances des modèles sur la prédiction de la pauvreté et de la vulnérabilité}
\label{tab:model_perf}
\begin{tabular}{llcccc}
\toprule
\textbf{Modèle} & \textbf{Cible} & \textbf{AUROC} & \textbf{Précision} & \textbf{Rappel} & \textbf{F1} \\
\midrule
LightGBM Baseline     & Pauvre     & 0.94 & 0.91 & 0.88 & 0.89 \\
LightGBM Baseline     & Vulnérable & 0.89 & 0.84 & 0.79 & 0.82 \\
LightGBM Transfert    & Pauvre     & 0.93 & 0.89 & 0.86 & 0.87 \\
LightGBM Transfert    & Vulnérable & 0.87 & 0.82 & 0.76 & 0.79 \\
Hyperréseau PyTorch   & Pauvre     & 0.91 & 0.86 & 0.83 & 0.85 \\
Hyperréseau PyTorch   & Vulnérable & 0.85 & 0.79 & 0.74 & 0.76 \\
\bottomrule
\end{tabular}
\end{table}

\includegraphics[width=12\linewidth,height=6.5cm,keepaspectratio]{compare.png}

\section{Interprétation des résultats}

\subsection{Pourquoi LightGBM domine-t-il~?}

Le LightGBM baseline surpasse l'hyperréseau sur toutes les métriques. Ce n'est pas un échec~: c'est un résultat informatif qui s'explique par deux mécanismes~:

\begin{enumerate}
  \item \textbf{Goulot d'embedding}~: Comprimer la région et le milieu dans un vecteur de dimension 16 perd de l'information que LightGBM capture directement depuis les catégories encodées et leurs interactions.
  \item \textbf{Déséquilibre de représentation}~: Les couches d'embedding ont besoin d'un volume suffisant d'exemples par catégorie pour converger. Certaines régions rurales peu représentées dans l'ENCDM ne fournissent pas assez de signal pour que leurs embeddings soient représentatifs.
\end{enumerate}

\subsection{La valeur de l'apprentissage par transfert}

Le LightGBM par transfert, entraîné sur les embeddings RGPH puis ajusté sur l'ENCDM, atteint un AUROC de 0.93 -- seulement 0.01 en dessous du baseline entraîné directement sur l'ENCDM. Cette proximité est une découverte importante~: \textbf{les indicateurs structurels du recensement permettent d'estimer la pauvreté avec une précision proche de 99\,\% de celle obtenue avec les données de dépenses directes}. Dans des contextes où l'enquête de consommation est indisponible (régions non couvertes, nouvelles zones d'intérêt), les données du recensement seul peuvent servir de base de prédiction robuste.

\subsection{Importances des features et validation analytique}

\begin{table}[H]
\centering
\small
\caption{Top 5 features -- classifieur de pauvreté LightGBM Baseline}
\label{tab:feature_importance}
\begin{tabularx}{\textwidth}{@{}c l c >{\raggedright\arraybackslash}X@{}}
\toprule
\textbf{Rang} & \textbf{Feature} & \textbf{Score} & \textbf{Interprétation} \\
\midrule
1 & Poids de sondage (\texttt{coef\_menage}) & 0.18 & Signal socio-économique au niveau de la strate \\
2 & Dépense annuelle par personne (\texttt{DAP}) & 0.14 & Mesure directe du niveau de vie \\
3 & Niveau d'éducation du chef & 0.12 & Effet capital humain \\
4 & Taille du ménage & 0.09 & Économies d'échelle et charge de dépendance \\
5 & Âge du chef de ménage & 0.07 & Effets de cycle de vie \\
\bottomrule
\end{tabularx}
\end{table}

La dominance de l'éducation (rang 3) dans les importances LGBM confirme l'intuition analytique du chapitre précédent~: le gradient éducatif identifié dans les graphiques descriptifs est bien capturé par le modèle comme signal prédictif de premier ordre, et non comme artefact de corrélation.

\vspace{1cm}
\includegraphics[width=5\linewidth,height=4.5cm,keepaspectratio]{newplot (4).png}
\includegraphics[width=5\linewidth,height=4.5cm,keepaspectratio]{newplot (5).png}

% ============================================================
%  CHAPITRE 7 — TABLEAU DE BORD
% ============================================================
\chapter{Déploiement du tableau de bord Streamlit}

\section{Philosophie de conception}

Le tableau de bord a été conçu comme un \textbf{instrument de recherche institutionnel}, pas comme une vitrine de présentation. Les choix de design en découlent~:

\begin{itemize}
  \item \textbf{Esthétique institutionnelle}~: fond papier chaud, texte charcoal, palette restreinte (bleu, teal, ambre) pour une lisibilité prolongée.
  \item \textbf{Architecture modulaire}~: chaque onglet correspond à un module Python indépendant, permettant de les développer, tester et modifier séparément.
  \item \textbf{Zéro friction utilisateur}~: l'utilisateur ne voit jamais de codes entiers bruts, de valeurs standardisées, ni de messages d'erreur techniques. Toute la complexité de la pipeline est masquée derrière une interface lisible en français.
\end{itemize}

\section{Architecture modulaire du dashboard}

\begin{table}[H]
\centering
\small
\caption{Modules du tableau de bord et leurs responsabilités}
\label{tab:dashboard_modules}
\begin{tabularx}{\textwidth}{@{}L{3.8cm} >{\raggedright\arraybackslash}X@{}}
\toprule
\textbf{Module} & \textbf{Responsabilité} \\
\midrule
\texttt{data\_loader.py} & Chargement en cache des DataFrames, modèles et scalers \\
\texttt{translations.py} & Noms et codes de régions, mappings GeoJSON \\
\texttt{plots.py} & 22 visualisations Plotly avec palette unifiée \\
\texttt{mapping.py} & Carte choroplèthe et profils régionaux synchronisés \\
\texttt{network.py} & Graphes de dépendances d'imputation \\
\texttt{sandbox.py} & Formulaire de prédiction et inférence duale \\
\texttt{hypernet.py} & Reconstruction PyTorch depuis \texttt{Hypernet.pt} \\
\texttt{data\_preview.py} & Explorateurs paginés des tables CleanENCDM/RGPH \\
\texttt{theme.py} & Système CSS, typographie et composants \\
\texttt{components.py} & Primitives de mise en page (KPI, cartes, grilles) \\
\bottomrule
\end{tabularx}
\end{table}

\section{Onglet 1 --- Intégrité des données}

Ce premier onglet donne accès aux explorateurs de données paginés et filtrables pour CleanENCDM et CleanRGPH. Il répond à une exigence de transparence~: tout utilisateur du dashboard doit pouvoir vérifier les données sur lesquelles reposent les analyses, inspecter les valeurs imputées, et consulter les métadonnées de chaque colonne via les fichiers de mapping JSON.

\figplaceholder[Dashboard -- Intégrité des données]{dashboard\_tab1\_data.png}{Onglet ``Intégrité des données'' -- explorateur de la table CleanENCDM avec filtres par région et statut de pauvreté}{fig:dash_tab1}

\section{Onglet 2 --- Analyse régionale}

L'onglet central est le cœur analytique du dashboard. Lors de la sélection d'une région (par défaut~: Guelmim-Oued Noun), quatre graphiques Plotly se synchronisent immédiatement~:

\begin{enumerate}
  \item \textbf{Distribution des dépenses}~: histogramme pondéré des dépenses annuelles, superposé par statut de pauvreté (ménages pauvres en rouge, non-pauvres en bleu).
  \item \textbf{Accès aux équipements}~: barres horizontales RGPH -- pourcentages de ménages disposant de l'électricité, de l'eau, d'internet, d'un mobile, d'un véhicule.
  \item \textbf{Éducation et pauvreté}~: barres groupées des taux de pauvreté par niveau d'éducation, rendant le gradient visible directement.
  \item \textbf{Genre et pauvreté}~: donut chart et barres comparant les taux par genre du chef de ménage.
\end{enumerate}

\figplaceholder[Dashboard -- Analyse régionale]{dashboard\_tab2\_regional.png}{Onglet ``Analyse régionale'' -- vue 4 graphiques synchronisés pour la région Guelmim-Oued Noun (distribution des dépenses, équipements RGPH, gradient éducatif, genre)}{fig:dash_tab2}

\section{Onglet 3 --- Bac à sable prédictif}

Le troisième onglet est le module d'inférence interactive. L'utilisateur saisit les caractéristiques d'un ménage via des menus déroulants entièrement en français. En coulisses, la pipeline complète s'exécute~:

\begin{enumerate}
  \item Les entrées brutes sont mappées vers les codes entiers via les dictionnaires de mapping JSON.
  \item Les features manquantes sont imputées avec les imputers KNN/LGBM persistés.
  \item Les valeurs sont standardisées avec les StandardScalers persistés.
  \item L'inférence parallèle s'exécute simultanément sur LightGBM et le hyperréseau.
  \item Quatre cartes métriques affichent~: probabilité de pauvreté LGBM, probabilité de vulnérabilité LGBM, probabilité de pauvreté hyperréseau, probabilité de vulnérabilité hyperréseau.
\end{enumerate}

\figplaceholder[Dashboard -- Bac à sable]{dashboard\_tab3\_sandbox.png}{Bac à sable prédictif -- formulaire de saisie, 4 cartes métriques de résultats LightGBM / Hyperréseau, et graphique SHAP explicatif}{fig:dash_tab3}

\subsection{Le bouton ``Simuler le transfert rural''}

Ce bouton est le \textbf{moment signature} du dashboard. Il conserve toutes les caractéristiques du ménage inchangées et bascule uniquement le milieu de résidence d'Urbain à Rural. On observe alors~:
\begin{itemize}
  \item \textbf{LightGBM}~: la probabilité reste inchangée (modèle statique, insensible au contexte géographique dynamique).
  \item \textbf{Hyperréseau}~: la probabilité se modifie (il régénère ses poids depuis l'embedding rural mis à jour).
\end{itemize}

Cette divergence visualise concrètement la différence entre une inférence conditionnée au contexte et une inférence statique -- et illustre, par l'exemple, la valeur et les limites de chaque architecture.

\section{Performance et caching}

Toutes les ressources lourdes (DataFrames, modèles, figures) sont mises en cache avec \texttt{@st.cache\_resource}. Le premier chargement prend environ 8 secondes (lecture des deux fichiers Parquet, chargement des modèles)~; toutes les navigations ultérieures entre onglets et changements de région s'exécutent en moins d'une seconde.

% ============================================================
%  CHAPITRE 8 — SIMPUTE
% ============================================================
\chapter{Simpute : de l'internship à l'open source}

\section{Genèse et motivation}

\textbf{Simpute} est la contribution open source directement née du stage. Il ne s'agit pas d'un projet parallèle~: c'est la \textbf{généralisation} du moteur d'imputation développé pour la plateforme HCP, extrait de son couplage domaine-spécifique et reconstruit comme bibliothèque Python réutilisable et distribuable.

La plateforme HCP avait démontré un principe algorithmique général~: les données manquantes doivent être traitées colonne par colonne, avec l'algorithme sélectionné selon le profil de données de chaque cible. Une fois ce principe identifié, l'abstraction en bibliothèque générique s'imposait.

\begin{infobox}[Étymologie]
\textbf{Simpute} = \textbf{Smart} + \textbf{Impute}. ``Smart'' parce que le moteur ne suppose pas qu'un modèle unique convient à toutes les colonnes~: il profile ses entrées, sélectionne adaptativement l'algorithme optimal, échoue gracieusement avec un fallback, et produit des résultats reproductibles sur n'importe quelle entrée tabulaire.
\end{infobox}

\section{Architecture du moteur}

\subsection{Profiling des colonnes (\texttt{simpute/utils.py})}

Avant toute sélection d'algorithme, le module de profiling inspecte chaque colonne selon trois axes~:
\begin{enumerate}
  \item \textbf{Type de données}~: distinction numérique (continu vs discret) vs catégoriel (nominal vs ordinal) par heuristiques de dtype et seuils de cardinalité.
  \item \textbf{Taux de missingness}~: calcul de \texttt{null\_count / total\_count}. Si ce ratio dépasse 70\,\%, un avertissement est émis.
  \item \textbf{Cardinalité et distribution}~: mesure de l'asymétrie (skewness) et du nombre de valeurs uniques pour guider le routage algorithmique.
\end{enumerate}

\subsection{Table de routage algorithmique (\texttt{simpute/models.py})}

\begin{table}[H]
\centering
\caption{Table de routage du moteur Simpute -- sélection d'algorithme par profil de colonne}
\label{tab:simpute_routing}
\begin{tabularx}{\textwidth}{L{5.5cm}X}
\toprule
\textbf{Profil de colonne} & \textbf{Algorithme sélectionné} \\
\midrule
Catégorielle haute cardinalité ($>$15 valeurs) & CatBoostClassifier / LightGBMClassifier \\
Catégorielle binaire ou faible cardinalité & LogisticRegression (L2) / LinearSVC \\
Continue asymétrique (skewness $>$ 1) & LightGBMRegressor / ExtraTreesRegressor \\
Continue normale ou uniforme & KNNRegressor / BayesianRidge \\
Taux de missingness $>$ 70\,\% & Fallback médiane / mode + avertissement \\
\bottomrule
\end{tabularx}
\end{table}

La sélection finale parmi les candidats est effectuée par \textbf{validation croisée sur un sous-échantillon} (15\,000 lignes maximum), évitant ainsi le choix purement heuristique qui produisait les premières implémentations aux résultats médiocres (R\textsuperscript{2} proche de 0 pour les colonnes numériques).

\subsection{Estimateur principal (\texttt{simpute/core.py})}

\begin{lstlisting}[caption={Interface publique de Simpute -- API compatible scikit-learn}]
import pandas as pd
from simpute import Simpute

df = pd.read_csv("data.csv")

imputer = Simpute(exclude=["id"])
filled = imputer.fit_transform(df)

new_data = pd.read_csv("new_data.csv")
filled_new = imputer.transform(new_data)
\end{lstlisting}

En interne, \texttt{fit\_transform} orchestre une boucle multi-passes~:
\begin{enumerate}
  \item Les colonnes numériques sont traitées en premier (données plus stables, moins ambiguës).
  \item À chaque passe, les colonnes déjà imputées servent de prédicteurs pour les colonnes suivantes.
  \item Jusqu'à trois passes sont exécutées~: à chaque passe, des imputers plus précis (bénéficiant de voisins mieux estimés) peuvent affiner les imputations des passes précédentes.
  \item Si un modèle échoue à converger, un fallback médiane/mode est appliqué pour cette colonne.
\end{enumerate}

\section{Validation et assurance qualité}

\subsection{Garanties contractuelles (\texttt{tests/guard.py})}

Le module de tests encode les \textbf{invariants contractuels} de la bibliothèque~:
\begin{itemize}
  \item \textbf{Complétude}~: aucun \texttt{NaN} dans la sortie.
  \item \textbf{Validité des plages}~: toutes les valeurs imputées dans les plages observées de chaque colonne.
  \item \textbf{Seuil de précision adaptatif}~: $\max(0.5,\; \text{taux classe majoritaire} - 0.05)$ pour les colonnes catégorielles.
  \item \textbf{Reproductibilité}~: deux appels avec la même graine produisent un output identique au bit près.
  \item \textbf{Ordre fit-transform}~: exception explicite si \texttt{transform} est appelé sans \texttt{fit} préalable.
\end{itemize}

\subsection{Stratégie de benchmarking (\texttt{tests/mask.py})}

Pour évaluer la précision de reconstruction, le module sélectionne aléatoirement 15\,\% des valeurs non nulles de chaque colonne, les remplace par \texttt{NaN}, et compare les prédictions aux valeurs originales après imputation~:

\begin{table}[H]
\centering
\caption{Résultats de benchmarking Simpute sur jeu de données de test synthétique}
\label{tab:simpute_results}
\begin{tabular}{llcc}
\toprule
\textbf{Colonne} & \textbf{Type} & \textbf{Métrique} & \textbf{Valeur} \\
\midrule
\texttt{Pre\_Semester\_GPA} & Continue & MAE & 0.393 \\
\texttt{Tool\_Diversity} & Continue & MAE & 0.762 \\
\texttt{Paid\_Subscription} & Nominale binaire & Exactitude & 60.4\,\% \\
\texttt{Major\_Category} & Nominale multi & Exactitude & 47.0\,\% \\
\texttt{Year\_of\_Study} & Ordinale & Exactitude & $\approx$20.0\,\% \\
\bottomrule
\end{tabular}
\end{table}

La faible précision de \texttt{Year\_of\_Study} est attendue et ne constitue pas un défaut~: dans ce jeu de données de test, l'année d'études ne porte presque aucun signal prédictif dans les autres features disponibles. Le moteur se comporte correctement -- il essaie d'apprendre une structure qui n'existe pas, échoue gracieusement, et revient au fallback sans planter l'exécution.

\section{Publication sur PyPI}

\subsection{Configuration du packaging}

Le package a été configuré pour la distribution open source via \texttt{pyproject.toml}~:

\begin{lstlisting}[caption={Configuration du packaging Simpute (pyproject.toml)}]
[project]
name = "simpute"
version = "1.0.0"
description = "Smart adaptive imputation engine for tabular data"
license = "MIT"
requires-python = ">=3.10"
authors = ["Hvllvix"]
classifiers = [
    "Programming Language :: Python :: 3",
    "License :: OSI Approved :: MIT License",
    "Operating System :: OS Independent",
    "Topic :: Scientific/Engineering :: Artificial Intelligence",
]
\end{lstlisting}

\subsection{Workflow GitHub Actions}

Le workflow de release (\texttt{.github/workflows/release.yml}) s'exécute automatiquement à la publication d'une GitHub Release. La séquence est~:
\begin{enumerate}
  \item Checkout du code source tagué.
  \item Installation des outils de build Python (\texttt{pip}, \texttt{build}).
  \item Construction de la distribution source et de la wheel (\texttt{python -m build}).
  \item Upload vers PyPI via \textit{trusted publishing} OIDC (sans token API en clair).
\end{enumerate}

\section{Organisation finale du dépôt Simpute}

\begin{lstlisting}[style=treenostyle, caption={Structure du depot Simpute}]
simpute/
|-- .github/workflows/release.yml
|-- simpute/
|   |-- __init__.py
|   |-- core.py
|   |-- models.py
|   |-- utils.py
|-- tests/
|   |-- __init__.py
|   |-- data/test.csv
|   |-- guard.py
|   |-- mask.py
|-- Assets/Plots/
|-- pyproject.toml
|-- setup.py
|-- README.md
|-- LICENSE
|-- .gitignore
\end{lstlisting}

% ============================================================
%  CHAPITRE 9 — RÉSULTATS ET DISCUSSION
% ============================================================
\chapter{Résultats, apports et discussion}

\section{Résultats obtenus}

\subsection{Résultats analytiques}

\begin{resultbox}
\textbf{Déterminants structurels de la pauvreté identifiés~:}\\[4pt]
$\bullet$\quad Guelmim-Oued Noun~: taux de pauvreté de $\approx$18\,\% --- le plus élevé parmi les régions étudiées\\
$\bullet$\quad Écart de dépenses urbain/rural~: \textbf{2.3x} (45\,000 vs 22\,000 MAD/an)\\
$\bullet$\quad Gradient éducatif~: chaque palier réduit le risque de pauvreté de 40 à 50\,\%\\
$\bullet$\quad Écart de genre persistant~: +3.6 points de \% pour les ménages dirigés par une femme\\
$\bullet$\quad Corrélation qualité du logement / pauvreté~: r = 0.71 (deux sources indépendantes)
\end{resultbox}

\subsection{Résultats de modélisation}

\begin{resultbox}
\textbf{Performances du meilleur modèle (LightGBM Baseline)~:}\\[4pt]
$\bullet$\quad AUROC pauvreté~: \textbf{0.94}\\
$\bullet$\quad F1-score pauvreté~: \textbf{0.89}\\
$\bullet$\quad AUROC vulnérabilité~: \textbf{0.89}\\
$\bullet$\quad F1-score vulnérabilité~: \textbf{0.82}\\
$\bullet$\quad Transfert RGPH $\to$ ENCDM~: AUROC = 0.93 (99\,\% de la performance baseline)
\end{resultbox}

\subsection{Résultats de la plateforme}

Un tableau de bord Streamlit complet a été déployé, comprenant 22 visualisations Plotly dynamiques, un moteur d'inférence duale interactive, un explorateur de données paginé, et une carte choroplèthe nationale. Le temps de chargement initial est de 8 secondes~; la navigation entre onglets est sous la seconde.

\subsection{Résultats Simpute}

Simpute v1.0.0 est disponible sur PyPI (\texttt{pip install simpute}), accompagné d'une suite de tests contractuels, d'une documentation complète et d'un workflow CI/CD entièrement automatisé.

\section{Apports techniques}

Ce projet a mobilisé et approfondi un spectre large de compétences techniques~:

\begin{itemize}
  \item \textbf{Ingénierie des données}~: lecture SPSS via \texttt{pyreadstat}, pipelines idempotents, persistance multi-format (Parquet, JSON, Joblib), validation aller-retour.
  \item \textbf{Machine Learning avancé}~: LightGBM avec poids de sondage, optimisation F1 sur classes déséquilibrées, validation croisée stratifiée, analyse SHAP.
  \item \textbf{Deep Learning}~: architecture PyTorch originale (hyperréseau à embeddings géographiques), entraînement CPU, sérialisation de checkpoints complets.
  \item \textbf{Développement logiciel}~: architecture modulaire stricte, design patterns (caching, séparation des couches, registre d'artefacts), packaging Python (PyPI, CI/CD).
  \item \textbf{Développement assisté par IA}~: utilisation productive de Cursor et Gemini pour accélérer les tâches répétitives tout en conservant la maîtrise des choix algorithmiques.
  \item \textbf{Visualisation}~: 22 graphiques Plotly interactifs, carte choroplèthe GeoJSON, graphiques SHAP, système de design CSS cohérent.
\end{itemize}

\section{Apports analytiques}

Le projet a produit une compréhension structurée des mécanismes de la pauvreté dans la région Guelmim-Oued Noun et au Maroc plus largement. Il démontre que les données statistiques officielles, correctement traitées et visualisées, peuvent transformer des fichiers SPSS bruts en informations directement actionnables pour les décideurs.

La corrélation RGPH-ENCDM (r = 0.71) est un résultat analytique en soi~: elle valide la cohérence des deux sources et ouvre la voie à l'utilisation des données de recensement comme proxy de pauvreté dans des zones non couvertes par les enquêtes de consommation.

\section{Difficultés rencontrées et solutions adoptées}

\begin{warningbox}
\small
\textbf{Difficultés techniques et solutions~:}\\[4pt]
$\bullet$\quad \textbf{Hétérogénéité SPSS}~: variantes de noms de colonnes $\to$ résolveur d'alias\\[3pt]
$\bullet$\quad \textbf{Scale RGPH}~: 730\,000 lignes $\times$ 24 cibles $\to$ hyperparamètres dynamiques\\[3pt]
$\bullet$\quad \textbf{Hyperréseau}~: convergence lente en CPU $\to$ embedding réduit, régularisation\\[3pt]
$\bullet$\quad \textbf{Simpute v0}~: R\textsuperscript{2} proche de 0 $\to$ validation croisée pour la sélection\\[3pt]
$\bullet$\quad \textbf{Listings LaTeX}~: accents dans le code $\to$ contenu ASCII et style sans \texttt{literate}
\end{warningbox}

\section{Discussion critique}

\textbf{La donnée brute n'est pas la connaissance.} Les fichiers ENCDM et RGPH existaient avant ce projet. Leur valeur analytique était largement inexploitée, dispersée dans des formats non interopérables. Le pipeline de prétraitement est la condition préalable à toute analyse~: c'est là que se joue l'essentiel de l'effort, et c'est là que la plupart des projets académiques échouent par impatience.

\textbf{Le contexte géographique améliore l'explication, pas la prédiction.} L'hyperréseau ne surpasse pas LightGBM en précision, mais il se comporte différemment lors de la simulation de transfert rural. Cette différence de comportement -- invisible dans les métriques AUROC -- est précisément ce qui intéresse les décideurs~: comprendre comment un même ménage serait affecté par un changement de contexte territorial.

\textbf{L'abstraction en open source est une compétence à part entière.} Extraire Simpute de la plateforme HCP a requis de distinguer ce qui appartient au domaine spécifique de la pauvreté marocaine de ce qui constitue un principe général d'imputation adaptative. Cette capacité d'abstraction -- reconnaître le général dans le particulier -- est l'une des compétences les plus précieuses en ingénierie logicielle.

% ============================================================
%  CONCLUSION GÉNÉRALE
% ============================================================
\chapter*{Conclusion générale et perspectives}
\addcontentsline{toc}{chapter}{Conclusion générale et perspectives}

\section*{Synthèse du projet}

Ce projet de fin d'études avait pour ambition de dépasser la simple analyse descriptive pour construire une véritable plateforme de data intelligence à partir des données officielles du Haut-Commissariat au Plan. Cet objectif a été atteint à plusieurs niveaux~:

Sur le plan \textbf{technique}, une chaîne analytique complète a été construite~: de l'ingestion de fichiers SPSS téléchargés depuis \texttt{hcp.ma} jusqu'à un tableau de bord Streamlit déployé, en passant par un pipeline de prétraitement robuste, des modèles d'apprentissage automatique et deep learning, et une bibliothèque open source publiée sur PyPI. L'utilisation de \textbf{Cursor} et \textbf{Gemini} comme outils d'assistance IA a significativement accéléré les phases de codage et de recherche documentaire.

Sur le plan \textbf{analytique}, les données ont révélé que la pauvreté au Maroc est un phénomène structurel ancré dans des inégalités d'infrastructure, d'éducation et de genre. Les modèles confirment ces patterns~: LightGBM atteint un AUROC de 0.94, et l'apprentissage par transfert démontre que les indicateurs de recensement seuls permettent de prédire la pauvreté à 99\,\% de la précision obtenue avec les données de dépenses directes.

Sur le plan \textbf{professionnel}, ce stage a offert une immersion concrète dans le fonctionnement d'une institution statistique nationale, avec ses contraintes de format, ses exigences de rigueur méthodologique et ses enjeux de valorisation de l'information au service du territoire.

\section*{Perspectives d'amélioration}

\begin{enumerate}
  \item \textbf{Connecter le dashboard à des données dynamiques}~: intégrer une base de données relationnelle permettant des mises à jour automatiques à chaque nouvelle vague d'enquête HCP.
  \item \textbf{Déploiement web public}~: migrer de Streamlit local vers une infrastructure cloud (FastAPI + frontend React) pour une diffusion externe accessible aux acteurs du développement régional.
  \item \textbf{Enrichissement du hyperréseau}~: concaténer les embeddings géographiques aux features brutes avant la couche finale -- une architecture hybride non testée dans ce projet qui pourrait améliorer les performances.
  \item \textbf{Simpute v2.0}~: ajouter le support des séries temporelles, des matrices creuses et des intervalles de confiance sur les imputations.
  \item \textbf{Extension régionale}~: adapter la plateforme aux onze autres directions régionales du HCP pour une couverture nationale, avec mise à jour automatique à chaque publication RGPH.
  \item \textbf{Intégration RGPH 2024}~: le recensement de 2024 est en cours de traitement au HCP. Ses données constitueront une mise à jour majeure pour la plateforme.
\end{enumerate}

\vspace{1cm}
\begin{flushright}
\begin{minipage}{0.6\textwidth}
\itshape ``Les données manquantes ne sont pas un défaut à minimiser~; elles sont une caractéristique structurelle des données du monde réel qui exige une réponse algorithmique réfléchie.''\\[6pt]
{\normalsize\upshape --- Réflexion issue du développement de Simpute}
\end{minipage}
\end{flushright}

% ============================================================
%  BIBLIOGRAPHIE
% ============================================================
\newpage
\chapter*{Références bibliographiques}
\addcontentsline{toc}{chapter}{Références bibliographiques}

\begin{enumerate}[label={[\arabic*]}, leftmargin=1.8cm]
  \item Haut-Commissariat au Plan, \textit{Enquête Nationale sur la Consommation et les Dépenses des Ménages 2019--2020 -- Rapport principal}, HCP, Rabat, 2021. Disponible sur~: \url{https://www.hcp.ma}
  \item Haut-Commissariat au Plan, \textit{Recensement Général de la Population et de l'Habitat 2014 -- Résultats détaillés}, HCP, Rabat, 2015. Disponible sur~: \url{https://www.hcp.ma}
  \item Ke, G., Meng, Q., Finley, T., Wang, T., Chen, W., Ma, W., Ye, Q., Liu, T.-Y., \textit{LightGBM: A Highly Efficient Gradient Boosting Decision Tree}, Advances in Neural Information Processing Systems (NeurIPS), 2017.
  \item Ha, D., Dai, A., Le, Q. V., \textit{HyperNetworks}, International Conference on Learning Representations (ICLR), 2017.
  \item Lundberg, S. M., Lee, S.-I., \textit{A Unified Approach to Interpreting Model Predictions (SHAP)}, NeurIPS, 2017.
  \item Troyanskaya, O., et al., \textit{Missing value estimation methods for DNA microarrays (KNN Imputation)}, Bioinformatics, 2001.
  \item Documentation officielle Python~: \url{https://www.python.org}
  \item Documentation Pandas~: \url{https://pandas.pydata.org}
  \item Documentation LightGBM~: \url{https://lightgbm.readthedocs.io}
  \item Documentation PyTorch~: \url{https://pytorch.org/docs}
  \item Documentation Streamlit~: \url{https://docs.streamlit.io}
  \item Documentation Plotly~: \url{https://plotly.com/python}
  \item Cursor IDE~: \url{https://cursor.sh}
  \item Google Gemini~: \url{https://gemini.google.com}
  \item Dépôt principal HCP-Internship (GitHub)~: \url{https://github.com/Hvllvix/HCP-Internship}
  \item Bibliothèque Simpute (PyPI)~: \url{https://pypi.org/project/simpute}
  \item Bibliothèque Simpute (GitHub)~: \url{https://github.com/Hvllvix/Simpute}
  \item Profil GitHub de l'auteur~: \url{https://github.com/Hvllvix}
\end{enumerate}

% ============================================================
%  ANNEXES
% ============================================================
\appendix

\chapter{Structure complète des dépôts GitHub}

\section{Dépôt HCP-Internship (branche Dashboard)}

\begin{lstlisting}[style=treenostyle, caption={Arborescence complete du depot HCP-Internship}]
HCP-Internship/
|-- Notebooks/
|   |-- Pre-processing.ipynb
|   |-- Analysis.ipynb
|   |-- Modeling.ipynb
|-- Dashboard/
|   |-- app.py
|   |-- data_loader.py
|   |-- translations.py
|   |-- plots.py
|   |-- mapping.py
|   |-- network.py
|   |-- sandbox.py
|   |-- hypernet.py
|   |-- data_preview.py
|   |-- theme.py
|   |-- components.py
|-- Artifacts/
|   |-- scalers/
|   |-- imputers/
|   |-- mappings/
|   |-- models/
|       |-- LGBMPauvre.pkl
|       |-- LGBMVulnerable.pkl
|       |-- Hypernet.pt
|-- Data/
|   |-- CleanENCDM.parquet
|   |-- CleanRGPH.parquet
|-- requirements.txt
|-- README.md
|-- .gitignore
\end{lstlisting}

\section{Dépôt Simpute}

\begin{lstlisting}[style=treenostyle, caption={Arborescence complete du depot Simpute}]
simpute/
|-- .github/
|   |-- workflows/
|       |-- release.yml
|-- simpute/
|   |-- __init__.py
|   |-- core.py
|   |-- models.py
|   |-- utils.py
|-- tests/
|   |-- __init__.py
|   |-- data/
|   |   |-- test.csv
|   |   |-- README.md
|   |-- guard.py
|   |-- mask.py
|-- Assets/
|   |-- Plots/
|-- pyproject.toml
|-- setup.py
|-- README.md
|-- LICENSE
|-- .gitignore
\end{lstlisting}

\chapter{Extrait du notebook de modélisation}

\section{Entraînement LightGBM avec poids de sondage}

\begin{lstlisting}[caption={Entrainement du classifieur LightGBM avec poids de sondage ENCDM}]
import joblib
import numpy as np
from lightgbm import LGBMClassifier
from sklearn.model_selection import StratifiedKFold, cross_val_score

df = pd.read_parquet("Artifacts/Data/CleanENCDM.parquet")

target = "Pauvre"
exclude = ["Pauvre", "Vulnerable", "coef_menage", "coef_indiv"]
features = [c for c in df.columns if c not in exclude]

X = df[features].values
y = df[target].values
w = df["coef_menage"].values

lgbm = LGBMClassifier(
    n_estimators=500,
    num_leaves=63,
    learning_rate=0.05,
    min_child_samples=10,
    class_weight="balanced",
    random_state=42,
)

cv = StratifiedKFold(n_splits=5, shuffle=True, random_state=42)
scores = cross_val_score(
    lgbm, X, y, cv=cv, scoring="roc_auc",
    sample_weight=w,
)

print("AUROC moyen : %.4f +/- %.4f" % (scores.mean(), scores.std()))

lgbm.fit(X, y, sample_weight=w)
joblib.dump(lgbm, "Artifacts/models/LGBMPauvre.pkl")
\end{lstlisting}

\end{document}
